
\chapter{Rokhsar--Kivelson Points and Quantum Dimer Models}
\label{chap:rk-qdm}

The toric code and the quantum double models are exactly solvable because their
Hamiltonians are sums of mutually commuting projectors.  The Kitaev honeycomb
model is solvable because the spin Hamiltonian reduces to free Majorana
fermions in every static gauge sector.  Rokhsar--Kivelson (RK) quantum dimer
models provide a third mechanism in two dimensions: the full Hamiltonian need
not be diagonalizable, but a special point in parameter space has an exact
frustration-free ground state whose amplitudes are controlled by a classical
statistical model.

The central idea is
\boxedpar{%
classical constrained configurations
\(\mathcal{C}\)
\(\Longrightarrow\)
quantum Hilbert space \(\operatorname{span}\{\ket{C}:C\in\mathcal{C}\}\)
\(\Longrightarrow\)
local resonance Hamiltonian
\(\Longrightarrow\)
exact equal-amplitude or weighted ground state at the RK point.}
This chapter explains this mechanism for quantum dimer models.  The model is
not usually exactly solvable in the strong spectral sense.  Its exact content is
more specific: the ground-state wavefunction, equal-time ground-state
correlators, and topological-sector structure can be obtained from a classical
dimer ensemble.

\section{Classical dimer coverings and the quantum Hilbert space}
\label{sec:qdm-hilbert-space}

Let \(\Lambda\) be a two-dimensional lattice with vertex set \(V\) and edge set
\(E\).  A \emph{dimer covering}, or perfect matching, is a subset
\(D\subset E\) such that every vertex is incident on exactly one occupied edge.
Equivalently,
\begin{equation}
\label{eq:perfect-matching-constraint}
    \sum_{\ell\ni i} n_{\ell}=1,
    \qquad
    n_{\ell}\in\{0,1\},
    \qquad
    i\in V .
\end{equation}
The classical configuration space is therefore
\begin{equation}
\label{eq:dimer-configuration-space}
    \mathcal{C}_{\rm dim}
    =
    \{D\subset E:\text{\(D\) satisfies \cref{eq:perfect-matching-constraint}}\}.
\end{equation}
The quantum dimer Hilbert space is the vector space spanned by orthonormal basis
states labelled by dimer coverings:
\begin{equation}
\label{eq:qdm-hilbert-space}
    \calH_{\rm QDM}
    =
    \operatorname{span}\{\ket{D}:D\in\mathcal{C}_{\rm dim}\},
    \qquad
    \braket{D}{D'}=\delta_{D,D'} .
\end{equation}
This Hilbert space is already constrained.  It is not the full tensor product of
independent qubits on edges.  The local condition \cref{eq:perfect-matching-constraint}
plays a role analogous to a Gauss-law constraint: every site must be touched by
exactly one dimer.

\begin{remark}[Dimer constraint as an emergent gauge constraint]
On a bipartite lattice one can orient every edge from the \(A\) sublattice to
the \(B\) sublattice and define a background charge.  The dimer constraint can
then be rewritten as a lattice divergence equation.  This is why height fields,
Coulomb phases, and gauge-theory language naturally appear in quantum dimer
models.  The constraint is not optional; it defines the Hilbert space.
\end{remark}

The elementary local move in the square-lattice quantum dimer model is a
plaquette flip.  A plaquette \(p\) is called \emph{flippable} if it contains two
parallel dimers on opposite edges.  There are two such local states, denoted
abstractly by
\begin{equation}
\label{eq:local-dimer-plaquette-states}
    \ket{h_p},
    \qquad
    \ket{v_p},
\end{equation}
where \(h\) and \(v\) stand for the two parallel dimer orientations on the
plaquette.  The plaquette flip exchanges these two states and leaves all other
edges unchanged.

\section{The square-lattice quantum dimer Hamiltonian}
\label{sec:square-lattice-qdm-hamiltonian}

The standard square-lattice Rokhsar--Kivelson quantum dimer Hamiltonian is
\begin{equation}
\label{eq:qdm-hamiltonian-square}
    H_{\rm QDM}
    =
    \sum_p H_p,
\end{equation}
with local plaquette term
\begin{equation}
\label{eq:qdm-local-hamiltonian}
    H_p
    =
    -t
    \left(
        \ket{h_p}\bra{v_p}
        +
        \ket{v_p}\bra{h_p}
    \right)
    +
    v
    \left(
        \ket{h_p}\bra{h_p}
        +
        \ket{v_p}\bra{v_p}
    \right).
\end{equation}
The first term is kinetic: it resonates two dimer coverings that differ by one
plaquette flip.  The second term is potential: it counts flippable plaquettes.
The parameters \(t\) and \(v\) are usually taken real, and we will assume
\(t>0\).

In the two-dimensional subspace
\(\operatorname{span}\{\ket{h_p},\ket{v_p}\}\), the local operator is the matrix
\begin{equation}
\label{eq:qdm-local-matrix}
    H_p
    \widehat{=}
    \begin{pmatrix}
        v & -t \\
        -t & v
    \end{pmatrix}.
\end{equation}
Thus the symmetric and antisymmetric local combinations have energies
\begin{equation}
\label{eq:qdm-local-eigenvalues}
    E_{+}=v-t,
    \qquad
    E_{-}=v+t,
\end{equation}
where
\begin{equation}
\label{eq:qdm-local-symmetric-antisymmetric}
    \ket{+_p}=\frac{1}{\sqrt{2}}(\ket{h_p}+\ket{v_p}),
    \qquad
    \ket{-_p}=\frac{1}{\sqrt{2}}(\ket{h_p}-\ket{v_p}).
\end{equation}
The RK point is the special point
\begin{equation}
\label{eq:rk-point-condition}
    v=t .
\end{equation}
At this point, the local Hamiltonian becomes a positive semidefinite projector
up to a scalar factor:
\begin{equation}
\label{eq:qdm-local-rk-projector}
    H_p^{\rm RK}
    =
    t
    \left(
        \ket{h_p}-\ket{v_p}
    \right)
    \left(
        \bra{h_p}-\bra{v_p}
    \right).
\end{equation}
Therefore
\begin{equation}
\label{eq:qdm-rk-positive-semidefinite}
    H_{\rm QDM}^{\rm RK}\geq0 .
\end{equation}
This is the frustration-free structure hidden in the model.

\begin{remark}[Not a commuting-projector model]
Although each \(H_p^{\rm RK}\) is a local positive-semidefinite projector, the
different plaquette terms do not generally commute.  Thus the RK point is not a
toric-code-type commuting-projector point.  Exact solvability here means exact
ground-state solvability, not simultaneous diagonalization of the full spectrum.
\end{remark}

\section{The equal-amplitude RK ground state}
\label{sec:equal-amplitude-rk-ground-state}

The plaquette flip dynamics decomposes \(\mathcal{C}_{\rm dim}\) into connected
components.  On a torus these components are labelled, in part, by winding
numbers.  Let \(\mathcal{S}\subset\mathcal{C}_{\rm dim}\) be one connected
sector under local plaquette flips.  Define
\begin{equation}
\label{eq:rk-equal-amplitude-state}
    \ket{\Psi_{\mathcal{S}}}
    =
    \frac{1}{\sqrt{N_{\mathcal{S}}}}
    \sum_{D\in\mathcal{S}}\ket{D},
    \qquad
    N_{\mathcal{S}}=\abs{\mathcal{S}} .
\end{equation}
This is the equal-amplitude superposition of all dimer coverings in the sector
\(\mathcal{S}\).

\begin{proposition}[RK ground state]
At \(v=t\), the state \(\ket{\Psi_{\mathcal{S}}}\) is annihilated by every local
RK plaquette term:
\begin{equation}
\label{eq:rk-state-annihilated}
    H_p^{\rm RK}\ket{\Psi_{\mathcal{S}}}=0
    \qquad
    \text{for all }p .
\end{equation}
Consequently,
\begin{equation}
\label{eq:rk-state-ground-energy}
    H_{\rm QDM}^{\rm RK}\ket{\Psi_{\mathcal{S}}}=0 .
\end{equation}
\end{proposition}

\begin{proof}
Fix a plaquette \(p\).  The only basis states on which
\(H_p^{\rm RK}\) acts nontrivially are pairs of dimer coverings
\((D_h,D_v)\) that differ by flipping \(p\).  In the equal-amplitude state,
\(\ket{D_h}\) and \(\ket{D_v}\) occur with the same coefficient.  The local
operator \cref{eq:qdm-local-rk-projector} projects onto the antisymmetric
combination \(\ket{D_h}-\ket{D_v}\).  Therefore the contribution from each
flippable pair cancels.  Nonflippable configurations are annihilated because
neither \(\ket{h_p}\) nor \(\ket{v_p}\) is present locally.  Hence every
\(H_p^{\rm RK}\) annihilates the state.
\end{proof}

The proof also explains the logic of the RK construction.  The Hamiltonian is
built so that every allowed local move imposes an equality of amplitudes between
configurations connected by that move.  If the sector \(\mathcal{S}\) is
connected under the local moves, those local equality conditions force a unique
equal-amplitude state within that sector.

\begin{theorem}[Sector-wise RK ground space]
Assume a connected component \(\mathcal{S}\) of the dimer configuration graph is
connected by the plaquette flips appearing in the Hamiltonian.  At \(v=t>0\),
the zero-energy ground state in that component is the equal-amplitude state
\(\ket{\Psi_{\mathcal{S}}}\).  The full RK ground space is spanned by one such
state for each dynamically disconnected sector, together with possible isolated
frozen configurations that contain no flippable plaquettes.
\end{theorem}

\begin{remark}[Frozen states]
Some lattices or boundary conditions admit staggered dimer configurations with
no flippable plaquettes.  These states are also annihilated by
\(H_{\rm QDM}^{\rm RK}\), because every local term acts only on flippable
plaquettes.  They are exact zero-energy states, but they are not liquid-like RK
superpositions.  When discussing universal long-distance physics, one usually
focuses on the liquid sectors rather than these frozen sectors.
\end{remark}

\section{Topological sectors and winding numbers}
\label{sec:qdm-topological-sectors}

On a simply connected domain with fixed boundary conditions, the local plaquette
flips typically connect all dimer coverings compatible with the boundary.  On a
torus, however, local flips cannot change global winding numbers.

For a bipartite square lattice, choose a reference dimer configuration \(D_0\).
Given another dimer covering \(D\), the superposition of \(D\) and \(D_0\) forms
oriented transition loops.  The net number of transition-loop crossings through
two noncontractible cuts defines two winding numbers
\begin{equation}
\label{eq:qdm-winding-numbers}
    (W_x,W_y)\in\ZZ\times\ZZ .
\end{equation}
A local plaquette flip changes the dimer configuration only inside a
contractible region.  Therefore it cannot change \((W_x,W_y)\).

\begin{proposition}[Conservation of winding sectors]
For the square-lattice quantum dimer model on a torus,
\begin{equation}
\label{eq:qdm-winding-conservation}
    [H_{\rm QDM},W_x]=[H_{\rm QDM},W_y]=0,
\end{equation}
where \(W_x\) and \(W_y\) denote the winding-sector labels regarded as operators
that are diagonal in the dimer basis.
\end{proposition}

\begin{proof}
Each Hamiltonian term flips dimers around a single plaquette.  Such a move adds
or removes an elementary contractible loop in the transition-graph
representation.  A contractible modification has zero intersection number with
any noncontractible cut.  Hence the winding numbers are unchanged by every
local term.
\end{proof}

The existence of winding sectors is similar in spirit to the topological sectors
of the toric code, but the physical meaning is different.  In the toric code the
fourfold degeneracy on a torus is robust throughout a topologically ordered
phase.  In the square-lattice RK dimer model, the number of winding sectors can
grow with system size, and the RK point is critical rather than a stable gapped
\(\ZZ_2\) topological phase.

\section{Relation to classical dimer partition functions}
\label{sec:rk-classical-dimer-correspondence}

The RK wavefunction converts quantum equal-time expectation values into
classical statistical averages.  Let \(\widehat{O}\) be an operator diagonal in
the dimer basis:
\begin{equation}
\label{eq:diagonal-dimer-operator}
    \widehat{O}\ket{D}=O(D)\ket{D} .
\end{equation}
Then in the equal-amplitude RK state,
\begin{equation}
\label{eq:rk-diagonal-observable-classical-average}
    \bra{\Psi_{\mathcal{S}}}\widehat{O}\ket{\Psi_{\mathcal{S}}}
    =
    \frac{1}{N_{\mathcal{S}}}
    \sum_{D\in\mathcal{S}}O(D).
\end{equation}
Thus diagonal ground-state correlations are exactly the correlations of a
classical dimer model with uniform Boltzmann weight in the same sector.

More generally, one may introduce positive classical weights \(w(D)>0\) and the
wavefunction
\begin{equation}
\label{eq:weighted-rk-wavefunction}
    \ket{\Psi_w}
    =
    \frac{1}{\sqrt{Z_w}}
    \sum_{D\in\mathcal{C}_{\rm dim}}
    \sqrt{w(D)}\ket{D},
    \qquad
    Z_w=
    \sum_{D\in\mathcal{C}_{\rm dim}}w(D).
\end{equation}
Then diagonal observables satisfy
\begin{equation}
\label{eq:weighted-rk-classical-average}
    \bra{\Psi_w}\widehat{O}\ket{\Psi_w}
    =
    \frac{1}{Z_w}
    \sum_{D\in\mathcal{C}_{\rm dim}}w(D)O(D).
\end{equation}
This is the broader classical-to-quantum RK construction.

\begin{principle}[RK classical--quantum correspondence]
If a quantum ground state has amplitudes equal to the square roots of classical
Boltzmann weights, then the norm of the quantum state is the classical partition
function, and all equal-time diagonal quantum correlators are classical thermal
correlators.
\end{principle}

For the uniform dimer RK state,
\begin{equation}
\label{eq:rk-norm-classical-partition-function}
    \braket{\Psi}{\Psi}
    =
    \sum_{D\in\mathcal{C}_{\rm dim}}1
    =
    Z_{\rm dimer} .
\end{equation}
This relation is the reason why classical dimer techniques, especially
Kasteleyn orientations and Pfaffians, enter the quantum problem.

\section{Kasteleyn orientation and Pfaffian solvability}
\label{sec:kasteleyn-pfaffian-dimers}

Classical close-packed dimers on a planar graph are exactly solvable by the
Kasteleyn--Pfaffian method.  Assign an antisymmetric matrix \(K\) to the graph,
with entries
\begin{equation}
\label{eq:kasteleyn-matrix-definition}
    K_{ij}
    =
    \begin{cases}
        \phantom{-}\kappa_{ij}, & \text{if the oriented edge is }i\to j,\\
        -\kappa_{ij}, & \text{if the oriented edge is }j\to i,\\
        0, & \text{if }i\text{ and }j\text{ are not adjacent},
    \end{cases}
\end{equation}
where \(\kappa_{ij}>0\) is the dimer fugacity on edge \((ij)\).  A Kasteleyn
orientation is an orientation for which each face has an odd number of clockwise
oriented edges, with a standard modification for boundary and topology.

For a planar graph with a Kasteleyn orientation, the classical dimer partition
function is
\begin{equation}
\label{eq:dimer-partition-pfaffian}
    Z_{\rm dimer}=\abs{\Pf K} .
\end{equation}
Since
\begin{equation}
\label{eq:pfaffian-determinant-relation}
    (\Pf K)^2=\det K,
\end{equation}
one can compute the free energy from a free-fermion determinant.  Moreover,
dimer occupation correlations are obtained from entries of \(K^{-1}\).  For
example, for two edges \(e=(ij)\) and \(e'=(kl)\), connected dimer correlations
can be expressed by finite determinants built from matrix elements of
\(K^{-1}\).

\begin{remark}[Why Pfaffians matter for the quantum RK point]
The RK Hamiltonian itself is not a free-fermion Hamiltonian.  However, its
exact ground-state norm and diagonal equal-time correlators reduce to a
classical dimer problem.  If that classical dimer problem is Pfaffian-solvable,
then those quantum ground-state quantities are exactly computable by Pfaffian
methods.  This is a ground-state correspondence, not a full spectral
fermionization of the quantum Hamiltonian.
\end{remark}

On a torus, one needs four Pfaffians corresponding to different boundary signs
around the two noncontractible cycles.  Schematically,
\begin{equation}
\label{eq:torus-dimer-four-pfaffians}
    Z_{\rm dimer}^{\rm torus}
    =
    \frac{1}{2}
    \left|
        -\Pf K_{00}
        +\Pf K_{10}
        +\Pf K_{01}
        +\Pf K_{11}
    \right|,
\end{equation}
where the subscripts indicate twists along the two cycles.  The precise sign
convention depends on the chosen Kasteleyn orientation, but the need for four
spin structures is universal.

\section{Height representation and the square-lattice RK point}
\label{sec:qdm-height-representation}

On bipartite planar lattices, close-packed dimer configurations can be mapped to
integer-valued height configurations on the dual lattice.  The construction is
local: crossing an edge changes the height by an amount that depends on whether
the edge is occupied by a dimer and on the sublattice convention.  The exact
microscopic increments are conventional, but the long-distance structure is not.
The coarse-grained height field \(h(\bm r)\) has an effective classical action
of Gaussian form
\begin{equation}
\label{eq:classical-dimer-height-action}
    S_{\rm cl}[h]
    =
    \frac{K}{2}
    \int \dd^2 r\,
    (\nabla h)^2
    +\text{irrelevant or locking terms} .
\end{equation}
For the square-lattice uniform dimer model, this height theory is critical.
Consequently, the square-lattice RK ground state has algebraic equal-time dimer
correlations.

At the RK point, the quantum dynamics adds a kinetic term to the height field.
The long-wavelength quantum theory is the quantum Lifshitz theory
\begin{equation}
\label{eq:quantum-lifshitz-action}
    S_{\rm QL}[h]
    =
    \frac{1}{2}
    \int \dd\tau\,\dd^2 r
    \left[
        (\partial_{\tau}h)^2
        +
        \kappa^2(\nabla^2 h)^2
    \right] .
\end{equation}
The dispersion relation is therefore
\begin{equation}
\label{eq:quantum-lifshitz-dispersion}
    \omega(\bm q)\sim \kappa \abs{\bm q}^2,
\end{equation}
so the dynamic critical exponent is
\begin{equation}
\label{eq:rk-dynamical-exponent}
    z=2 .
\end{equation}

\begin{remark}[Why this is not a relativistic critical point]
In most quantum critical points discussed earlier in one-dimensional free
fermion chains, the low-energy theory has \(z=1\) and spacetime scaling is
relativistic.  The square-lattice RK point is different: its effective action
contains \((\nabla^2 h)^2\), not \((\nabla h)^2\), in the quantum potential
energy.  The resulting \(\omega\sim q^2\) dispersion is a hallmark of RK-type
criticality.
\end{remark}

The square-lattice phase diagram near the RK point is often summarized as
follows.  For \(v/t>1\), flippable plaquettes cost energy and staggered frozen
states dominate.  For sufficiently smaller \(v/t\), crystalline valence-bond
solid phases appear.  The RK point \(v=t\) is a fine-tuned critical point
separating the staggered regime from resonating dimer physics.

\section{Triangular lattice and gapped \texorpdfstring{\(\ZZ_2\)}{Z2} dimer liquid}
\label{sec:triangular-rk-liquid}

The physics of the RK point depends strongly on the lattice.  On the triangular
lattice, the elementary move again flips two parallel dimers around a rhombus,
but the lattice is nonbipartite.  The absence of a bipartite height field changes
the long-distance behavior.  At the triangular-lattice RK point, the equal-weight
dimer liquid is gapped and realizes \(\ZZ_2\) topological order rather than a
critical height phase.

The distinction may be summarized as
\begin{center}
\begin{tabular}{L{0.28\linewidth}L{0.30\linewidth}L{0.30\linewidth}}
\toprule
Lattice type & Classical dimer ensemble & RK quantum phase \\
\midrule
Bipartite square lattice & Height/Coulomb phase with algebraic correlations & Critical RK point with \(z=2\) \\
Nonbipartite triangular lattice & Short-ranged dimer correlations & Gapped \(\ZZ_2\) topological dimer liquid \\
\bottomrule
\end{tabular}
\end{center}

On a torus, a gapped \(\ZZ_2\) dimer liquid has four topological sectors,
analogous to the four ground states of the toric code.  A useful way to label
them is by the parity of the number of dimers crossing two noncontractible cuts:
\begin{equation}
\label{eq:z2-dimer-winding-parities}
    (P_x,P_y)\in\ZZ_2\times\ZZ_2 .
\end{equation}
Local plaquette flips preserve these parities.  In the thermodynamic limit, the
four corresponding RK ground states become locally indistinguishable if the
phase is gapped and topologically ordered.

\begin{remark}[Relation to the toric code]
The triangular-lattice RK dimer liquid and the toric code are not identical
microscopic Hamiltonians.  However, they lie in the same broad topological
universality class when the dimer liquid is gapped and has \(\ZZ_2\) topological
order.  The toric code is the fixed-point commuting-projector representative;
the quantum dimer model is a constrained Hilbert-space realization with
resonance dynamics.
\end{remark}

\section{General RK construction from stochastic dynamics}
\label{sec:general-rk-stochastic-construction}

The dimer RK point is a special case of a general construction.  Let
\(\mathcal{C}\) be a classical configuration space with positive weights
\(w(C)>0\).  Suppose there is a set of local reversible moves
\(C\leftrightarrow C'\).  Define local operators of the form
\begin{equation}
\label{eq:general-rk-local-projector}
    H_{C,C'}
    =
    \gamma_{C,C'}
    \left(
        \sqrt{\frac{w(C')}{w(C)}}\ket{C}
        -
        \sqrt{\frac{w(C)}{w(C')}}\ket{C'}
    \right)
    \left(
        \sqrt{\frac{w(C')}{w(C)}}\bra{C}
        -
        \sqrt{\frac{w(C)}{w(C')}}\bra{C'}
    \right),
\end{equation}
with \(\gamma_{C,C'}>0\).  Then the state \(\ket{\Psi_w}\) in
\cref{eq:weighted-rk-wavefunction} is annihilated by every such term.  Hence
\begin{equation}
\label{eq:general-rk-hamiltonian}
    H_{\rm RK}[w]
    =
    \sum_{(C,C')}H_{C,C'}
\end{equation}
is frustration-free and has \(\ket{\Psi_w}\) as an exact ground state in every
connected component of the move graph.

This construction is closely related to classical Markov dynamics satisfying
detailed balance.  A Markov generator \(M\) with stationary distribution
\begin{equation}
\label{eq:rk-stationary-distribution}
    \pi(C)=\frac{w(C)}{Z_w}
\end{equation}
can often be transformed by a similarity transformation into a Hermitian RK
Hamiltonian:
\begin{equation}
\label{eq:markov-to-rk-similarity}
    H_{\rm RK}
    =
    -\Pi^{-1/2}M\Pi^{1/2},
    \qquad
    \Pi_{C,C'}=\pi(C)\delta_{C,C'} .
\end{equation}
Consequently, relaxation modes of the classical stochastic process map to
excitation modes of the quantum RK Hamiltonian.

\begin{principle}[Stochastic-matrix form]
RK Hamiltonians often have a stochastic-matrix form: after a diagonal similarity
transformation, the quantum Hamiltonian becomes a classical Markov generator.
The exact ground state is the square root of the stationary distribution, while
low-energy quantum modes mirror slow classical relaxation modes.
\end{principle}

\section{Excitations: resonons, visons, and monomers}
\label{sec:qdm-excitations}

Although the RK point gives an exact ground state, its excitations are generally
more difficult than those of a free-fermion or commuting-projector model.
Nevertheless, several important excitation types can be identified.

\paragraph{Resonons.}
On the square-lattice RK point, low-energy collective modes of the height field
are sometimes called resonons.  Their energy vanishes as
\begin{equation}
\label{eq:resonon-dispersion}
    E(\bm q)\sim \abs{\bm q}^2,
\end{equation}
consistent with the quantum Lifshitz exponent \(z=2\).

\paragraph{Visons.}
In a gapped \(\ZZ_2\) dimer liquid, a vison is a \(\ZZ_2\) vortex excitation.
Operationally, a vison can be detected by a string operator that measures the
parity of dimers crossing a path on the dual lattice.  Its mutual statistics
with a monomer excitation is analogous to the mutual statistics of \(m\) and
\(e\) particles in the toric code.

\paragraph{Monomers.}
A monomer is a violation of the close-packing constraint: a site is left
unmatched or has the wrong dimer occupancy.  In a strict quantum dimer Hilbert
space, monomers are excluded.  They can be introduced by enlarging the Hilbert
space.  Their confinement or deconfinement diagnoses whether the dimer phase is
crystalline, critical, or topologically ordered.

\section{Comparison with earlier exact models}
\label{sec:qdm-comparison-with-earlier-models}

Quantum dimer models clarify an important distinction in these notes.  Exact
solvability can refer to the full spectrum, to a sector-wise reduction, or only
to the ground-state wavefunction and its equal-time correlations.  The RK point
belongs to the last category.

\begin{center}
\small
\begin{tabular}{L{0.23\linewidth}L{0.33\linewidth}L{0.31\linewidth}}
\toprule
Model class & Solvable structure & What is exactly controlled \\
\midrule
TFIM/XY/Kitaev chain & Jordan--Wigner plus BdG diagonalization & Full spectrum and Gaussian correlators \\
XXZ/XYZ chains & Yang--Baxter algebra and Bethe ansatz & Spectrum through Bethe equations \\
Toric code & Commuting stabilizers/projectors & Full spectrum and anyon algebra at the fixed point \\
Kitaev honeycomb & Static \(\ZZ_2\) gauge fields plus free Majoranas & Free-fermion spectrum in each flux sector \\
RK quantum dimer model & Frustration-free stochastic-matrix form & Exact ground state and classical equal-time correlators \\
\bottomrule
\end{tabular}
\end{center}

Thus the RK model is exactly solvable in a weaker but conceptually important
sense.  It shows that a nontrivial many-body wavefunction can be known exactly
even when the full many-body spectrum is not.

\section{Summary}
\label{sec:qdm-summary}

The essential points of this chapter are:
\begin{enumerate}[label=(\roman*)]
    \item A quantum dimer model is defined on the constrained Hilbert space of
    perfect matchings.
    \item The square-lattice QDM has kinetic plaquette flips and a potential
    term that counts flippable plaquettes.
    \item At the RK point \(v=t\), the Hamiltonian becomes a sum of positive
    semidefinite local terms.
    \item In every connected sector, the equal-amplitude superposition of dimer
    coverings is an exact zero-energy ground state.
    \item The norm and diagonal equal-time correlators of the RK state are
    those of a classical dimer model.
    \item Classical dimer Pfaffian methods therefore compute quantum
    ground-state correlations at the RK point.
    \item Bipartite square-lattice RK physics is critical and described by a
    height/quantum-Lifshitz theory with \(z=2\).
    \item Nonbipartite triangular-lattice RK physics can realize a gapped
    \(\ZZ_2\) topological dimer liquid.
\end{enumerate}

The next chapter moves from RK ground-state solvability to more general exact
ground-state constructions, including valence-bond-solid states, stabilizer
cluster states, and parent Hamiltonians of tensor-network states.
